\documentstyle[%


righttag%


,11pt%


,amssymb%


,russian%               remove in Eng.


,izvru%                 Set izven% in Eng.


% ,izvru-ks             for brief communications


]{amsart}





%       Math definitions





\newcommand{\col}{\operatorname{col}}


\newcommand{\tts}[1]{}





%\hoffset=-15mm%                  For Eng.


\setcounter{page}{75}


\god{2000}


\nomer{6~(457)} %                        Remove in Eng.


%\nomer{Vol.\,44, No.\,6}%               Set in Eng.


%\str{\pageref{begin}--\pageref{end}}%  Set in Eng.


\udk{517.929.4\,:\,519.21}





\author{\it Р.И.\,Кадиев}


% NB:   ^^ remove in Eng.


\title{Достаточные условия устойчивости по части переменных


линейных стохастических систем\\ с последействием}





\date{}


%\pagestyle{myheadings}%         Set in Eng.


\pagestyle{plain} %              Remove in Eng.


\begin{document}


%\thispagestyle{plain}%          Set in Eng.


\maketitle


\label{begin}





Пусть $(\Omega, {\cal F}, ({\cal F})_{t\ge 0}, P)$ --- 


стохастический базис ([1], с.\,9);


$\bigvee \limits_{s=0}^tg(s)$ --- полная вариация 


функции $g(s)$ на отрезке $[0, t]$; $D^n$ --- 


линейное пространство $n$-мерных предсказуемых ([1], с.\,13)


случайных процессов на $[0, +\infty [$, траектории которых почти наверно 


(п.\,н.) непрерывны справа и имеют предел слева; $k^n$ --- линейное 


пространство $n$-мерных ${\cal F}_0$-измеримых случайных величин;


$Z = \col (z^1,\dotsc, z^m)$ ecть $m$-мерный семимартингал ([1],


с.\,73); $L^n(Z)$ --- линейное пространство предсказуемых 


$n\times m$-матриц на $[0, +\infty [$, строки которых локально интегрируемы 


по семимартингалу $Z$; $\lambda :[0, +\infty [ \rightarrow R_+$ --- 


некоторая неубывающая функция; $\lambda ^*$ --- мера, порожденная функцией 


$\lambda$; $L^\lambda_q$ --- линейное пространство скалярных функций на 


$[0, +\infty [$, суммируемых со степенью $q$ при $1 \le q < \infty $ по 


мере $\lambda^*$ и ограниченных в существенном при $q = \infty$ 


по мере $\lambda^*$; $R^n$ --- линейное пространство $n$-мерных векторов с 


нормой $|\cdot|$; $\|\cdot\|$ --- норма $l\times n$-матрицы,


согласованная с нормой 


$|\cdot|$; $\|\cdot\|_X$ --- норма в нормированном пространстве $X$;


$\|x\|_t = \sup \limits_{0 \le s \le t}|x(s)|$, если $x \in R^n$;


$\|H\|_t = \sup \limits_{0 \le s \le t}\|H\|$, если $H$ --- 


$l\times n$-матрица;


$1 \le p < \infty $; $c_p$ --- положительное число, зависящее от $p$,


в неравенстве (9.48) из ([1], с.\,65). 





Для заданного числа $p$ и положительной скалярной функции 


$\gamma (t)$, $t \in [0, +\infty [$, в [2] введены нормированные 


пространства $M_p^\gamma$ $(M_p^1 = M_p)$, $k_p^n$. 


Для любых $x \in D^n$, $x = \col(x^1,\dotsc, x^n)$, $0 < k < n$,


введем следующие обозначения: $y = \col(x^1,\dotsc, x^k)$, 


$h = \col(x^{k+1},\dotsc, x^n)$. Тогда $x = \col(y,h)$ и 


$D^n = D^k\times D^{n-k}$. 





В дальнейшем предположим, что семимартингал $Z$ представ\'им в виде $Z=b+c$, 


где $b$ --- предсказуемый случайный процесс локально ограниченной вариации, а 


$c$ --- локально квадратично интегрируемый мартингал ([1], с.\,28). 


Кроме того, все компоненты процесса $b$ и взаимные характеристики


$\langle c^i,c^j\rangle$ 


([1], с.\,48) всех компонент мартингала $c$ будем предполагать 


абсолютно непрерывными относительно меры $\lambda^*$. Последнее означает, что 


$$


b^i = \int_0^{\displaystyle.}a^i d\lambda,


\quad \langle c^i,c^j\rangle = \int_0^{\displaystyle .}


A^{ij}d\lambda\quad (i, j = 1,\dotsc, m). 


$$





Пусть $a = \col (a^1,\dotsc, a^m)$; $A = [A^{ij}]$ --- $n\times m$-матрица;


$a^+ = \col(|a^1|,\dotsc, |a^m|)$; $A^+ = [|A^{ij}|]$; если $B$ --- 


$n\times n$-матрица, то 


$B^1$ --- $k\times k$-матрица, полученная из матрицы $B$ 


зачеркиванием $n-k$ последних столбцов и строк, 


$B^2$ --- $k\times (n-k)$-матрица, полученная из матрицы $B$ 


зачеркиванием $k$ первых столбцов и $n-k$ последних строк, 


$B^3$ --- $(n-k)\times k$-матрица, полученная из матрицы $B$ 


зачеркиванием первых $k$ строк и последних $n - k$ столбцов, 


$B^4$ --- $(n-k)\times (n-k)$-матрица, полученная из матрицы $B$ 


зачеркиванием $k$ первых строк и столбцов. 





Рассмотрим уравнение 


\begin{equation}


dx(t) = (Vx)(t)dZ(t),\quad t \ge 0,


\end{equation}


где $(Vx)(t) = \Big(\int \limits_0^td_s{\cal R}_1(t,s)x(s),\dotsc, 


\int \limits_0^td_s{\cal R}_m(t,s)x(s)\Big)$, 


${\cal R}_i(t,s) = \sum \limits_{j=0}^{m_i}A_{ij}(t)r_{ij}(t,s)$, 


$A_{ij}$ --- $n\times n$-матрица, элементами которой являются 


предсказуемые случайные процессы, $r_{ij}(t,s)$ --- скалярная 


функция, определенная в области $\{(t,s): 0 \le s \le t < +\infty \}$, 


$i = 1,\dotsc, m$, $j = 0,\dotsc, m_i$.





Уравнение (1) рассматривается в предположении, что 


$$


\int_0^t(|H_1a^+| + \|H_1A^+H_1^\top \|)d\lambda < \infty


$$ 


п.\,н. для любого $t \ge 0$, где $H_1 = (H_1^1,\dotsc, H_1^m)$, 


$H_1^i = \sum \limits_{j=0}^{m_i}\|A_{ij}\|


\bigvee \limits_{s=0}^{\displaystyle .}r_{ij}(\cdot,s)$ при $i=1,\dotsc, m$. 


Тогда оператор $V$ действует из пространства $D^n$ в пространство 


$L^n(Z)$ и через любое $x(0) \in k^n$ проходит единственное 


(с точностью до $P$-эквивалентности) 


решение $x(t)$ уравнения (1) [3]. Это решение


имеет представление $x(t) = X(t)x(0)$, $t \ge 0$, 


где $X$ --- фундаментальная матрица уравнения (1) [2]. 





Частными случаями уравнения (1) являются два вида систем.





a) Линейная стохастическая система с сосредоточенным запаздыванием


\begin{equation}


dx(t) = (V^1x)(t)dZ(t),\quad t \ge 0, \quad


x(s) = 0,\quad s < 0, 


\end{equation}


где $(V^1x)(t) =


\Big(\sum \limits_{j=0}^{m_1}\ov A_{1j}(t)x(h_{1j}(t)),\dotsc, 


\sum \limits_{j=0}^{m_m}\ov A_{mj}(t)x(h_{mj}(t))\Big)$, 


$\ov A_{ij}$ --- $n\times n$-матрица, элементами которой являются 


предсказуемые случайные процессы, $h_{ij}$ --- измеримая 


функция такая, что $h_{ij}(t) \le t$, $t \in [0, +\infty [$,


$\lambda^*$ почти всюду (п.\,в.) при 


$i = 1,\dotsc, m$, $j = 0,\dotsc, m_i$.





Уравнение (2) рассматривается в предположении, что 


$$


\int_0^t(|H_2a^+| + \|H_2A^+H_2^\top \|)d\lambda < \infty


$$


п.\,н. для любого $t \ge 0$, где $H_2 = (H_2^1,\dotsc, H_2^m)$, 


$H_2^i = \sum \limits_{j=0}^{m_i}\|\ov A_{ij}\|$ при $i = 1,\dotsc, m$. 


Тогда уравнение (2) является частным случаем уравнения 


(1) и ${\cal R}_i(t,s)=\sum\limits_{j=0}^{m_i}\ov A_{ij}(t)r_{ij}(t,s)$, 


где $r_{ij}$ --- характеристическая функция множества 


$\{(t,s): h_{ij}(t) \le s \le t\}$, определенная в области 


$\{(t,s): 0 \le s \le t < +\infty \}$ при $i = 1,\dotsc, m$, 


$j = 0,\dotsc, m_i$.





б) Линейная система ``обыкновенных'' стохастических 


дифференциальных уравнений


\begin{equation}


dx(t) = (V^2x)(t)dZ(t),\quad t \ge 0,


\end{equation}


где $(V^2x)(t) = (A_1(t)x(t),\dotsc, A_m(t)x(t))$, 


$A_i$ --- $n\times n$-матрица, элементами которой являются 


предсказуемые случайные процессы при $i = 1,\dotsc, m$.


Кроме того, 


$\int \limits_0^t(|H_3a^+| + \|H_3A^+H_3^\top \|)d\lambda < \infty $ 


п.\,н. для любого $t \ge 0$, где $H_3 = (H_3^1,\dotsc, H_3^m)$, 


$H_3^i = \|A_i\|$ при $i = 1,\dotsc, m$. 


В этом случае 


${\cal R}_i(t,s) = A_i(t)r(t,s)$, 


где $r$ --- характеристическая функция множества 


$\{(t,s): s = t\}$, определенная в области 


$\{(t,s): 0 \le s \le t < \infty \}$ при $i = 1,\dotsc, m$.





Пусть $x_i$, $i = 1,\dotsc, n$, --- столбцы матрицы $X$. 


Тогда через $Y$ обозначим $k\times n$-матрицу, столбцами 


которой являются $y_i$, $i = 1,\dotsc, n$. По теореме 1 


из [4] уравнение (1) $py$--устойчиво, асимптотически 


$py$--устойчиво, экспоненциально $py$--устойчиво ([4], определение 


3) тогда и только тогда, когда для любого 


$x_0 \in k_p^n$ имеем $Yx_0 \in M_p$, $Yx_0 \in M_p^\gamma $, 


где $\gamma \ge \delta > 0$, $\lim \limits_{t\to +\infty }\gamma (t) = 


+\infty $, $Yx_0 \in M_p^\gamma $, где $\gamma (t) = \exp \{\beta t\}$, 


а $\beta $ --- некоторое положительное число.


Поэтому в дальнейшем будем изучать более общий вид устойчивости 


уравнения (1), объединяющий все виды устойчивости. 


\medskip





{\bf Определение.} Уравнение (1) назовем $M_p^\gamma y$-устойчивым, 


если для любого $x_0 \in k_p^n$ имеем $Yx_0 \in M_p^\gamma$.


\medskip





Уравнение (1) эквивалентно системе вида


\begin{equation}


\begin{split}


dy(t) &= [(V_1y)(t) + (V_2h)(t)]dZ(t),\quad t \ge 0, \\


dh(t) &= [(V_3y)(t) + (V_4h)(t)]dZ(t),\quad t \ge 0, \\


\end{split}


\end{equation}


где $(V_lx)(t) = \Big(\int \limits_0^td_s{\cal R}_1^l(t,s),\dotsc, 


\int \limits_0^td_s{\cal R}_m^l(t,s)\Big)$, 


${\cal R}_i^l(t,s) = \sum \limits_{j=0}^{m_i}


A_{ij}^l(t)r_{ij}(t,s)$ при $l = \ov{1,4}$, $i = 1,\dotsc, m$. 





В силу того, что через любое $x(0) \in k^n$ проходит 


единственное решение $x(t)$ уравнения (1), каждое из уравнений 


системы (4) в отдельности будет иметь единственное решение при 


любых $y(0) \in k^k$, $h \in D^{n-k}$ и $h(0) \in k^{n-k}$, 


$y \in D^k$ соответственно. Тогда второе уравнение системы (4) 


эквивалентно уравнению $h(t) = H(t)h(0) + (K_1V_3y)(t)$, 


$t \ge 0$, где $H$ --- фундаментальная матрица, $K_1$ --- оператор 


Коши для второго уравнения этой системы [2].


Отсюда и из первого уравнения системы (4) получим 


\begin{equation}


dy(t) = [(V_5y)(t) + (V_2(Hh(0)))(t)]dZ(t),\quad t \ge 0, 


\end{equation}


где $V_5 = V_1 + V_2K_1V_3$, и уравнение (1) $M_p^\gamma y$-устойчиво 


тогда и только тогда, когда для любого $x(0) \in k_p^n$ решение 


уравнения (5) принадлежит пространству $M_p^\gamma$. 





В дальнейшем предположим, что для уравнений (1)--(3) оператор 


$V_4$ такой, что имеет место равенство 


$(K_1V_3y)(t) = H(t)\int \limits_0^t(H(s))^{-1}(V_3y)(s)dZ(s)$. 


Далее будем считать, что $\ov V_1$, $\ov V_2$, $\ov V_3$ --- 


операторы, определяемые равенствами $\ov V_1y = V_1y$, 


$\ov V_2y = V_2(Hy)$, $\ov V_3 = H^{-1}V_3y$ соответственно.





\begin{thm}


Пусть существуют предсказуемые 


случайные процессы $d_i$ и положительные числа $\ov c_i$ 


при $i = 1,\dotsc, 6$ такие, что для любых $y \in D^k$, 


$\phi \in D^{n-k}$ имеют место неравенства 


$|(\ov V_1y)(\cdot)a(\cdot)| \le d_1(\cdot)\|y\|_{(\cdot)}$,


$\|(\ov V_1y)(\cdot)A(\cdot)((\ov V_1y)(\cdot))^\top\|


\le d_2(\cdot)\|y\|_{(\cdot)}$,


$|(\ov V_2\phi )(\cdot)a(\cdot)| \le d_3(\cdot)\|\phi\|_{(\cdot)}$,


$\|(\ov V_2\phi )(\cdot)A(\cdot)((\ov V_2\phi )(\cdot))^\top\|


\le d_4(\cdot)\|\phi \|_{(\cdot)}$,


$|(\ov V_3y)(\cdot)a(\cdot)| \le d_5(\cdot)\|y\|_{(\cdot)}$,


$\|(\ov V_3y)(\cdot)A(\cdot)((\ov V_3y)(\cdot))^\top\|


\le d_6(\cdot)\|y\|_{(\cdot)}$, 


$\int \limits_0^\infty d_i(s)d\lambda (s) \le \ov c_i$ 


п.\,н. при $i = 1,\dotsc, 6$ и


$\int \limits_0^{\displaystyle .}(V_2(Hh(0)))(s)dZ(s) \in M_{2p}$ 


при любом $h(0) \in k_{2p}^{n-k}$. Тогда уравнение $(1)$ $M_{2p}y$-устойчиво.


\end{thm}





Пусть $\xi $ --- скалярная функция на $[0, +\infty [$, 


локально суммируемая по функции $\lambda $, а 


$\gamma (t) = \exp \Big\{-\int \limits_0^t\xi (s)d\lambda (s)\Big\}$. 





\begin{cor}


Пусть существуют предсказуемые 


случайные процессы $\wh d_i$ и положительные числа $\wh c_i$ 


при $i = 1,\dotsc, 6$ такие, что для любых $y \in D^k$, 


$\phi \in D^{n-k}$ имеют место неравенства


$|(1/\gamma (\cdot))(\ov V_1(\gamma y))(\cdot)a(\cdot) + \xi y| \le 


\wh d_1(\cdot)\|y\|_{(\cdot)}$,


$\|(1/\gamma (\cdot))^2(\ov V_1(\gamma y))(\cdot)A(\cdot)


((\ov V_1(\gamma y))(\cdot))^\top\| \le \wh d_2(\cdot)\|y\|_{(\cdot)}$,


$|(1/\gamma (\cdot))(\ov V_2(\gamma \phi ))(\cdot)a(\cdot)| \le 


\wh d_3(\cdot)\|\phi\|_{(\cdot)}$,


$\|(1/\gamma (\cdot))^2(\ov V_2(\gamma \phi ))(\cdot)A(\cdot)


((\ov V_2(\gamma \phi ))(\cdot))^\top\|


\le \wh d_4(\cdot)\|\phi \|_{(\cdot)}$,\linebreak


$|(1/\gamma (\cdot))(\ov V_3(\gamma y))(\cdot)a(\cdot)|


\le \wh d_5(\cdot)\|y\|_{(\cdot)}$,


$\|(1/\gamma (\cdot))^2(\ov V_3(\gamma y))(\cdot)A(\cdot)


((\ov V_3(\gamma y))(\cdot))^\top\|


\le \wh d_6(\cdot)\|y\|_{(\cdot)}$, \linebreak


$\int \limits_0^\infty \wh d_i(s)d\lambda (s) \le \wh c_i$ 


п.\,н. при $i = 1,\dotsc, 6$ и 


$\int \limits_0^{\displaystyle .}(1/\gamma(s))(V_2(Hh(0)))(s)dZ(s)\in M_{2p}$ 


при любом $h(0) \in k_{2p}^{n-k}$. Тогда уравнение (1) 


$M_{2p}^{1/\gamma }y$-устойчиво.


\end{cor}





\begin{cor}


Если существует $\wh C > 0$ 


такое, что для уравнения (1) имеет место неравенство 


$\int \limits_0^\infty (|H_1^ja^+| + 


\|H_1^jA^+(H_1^j)^\top \|)d\lambda \le \wh C$ п.\,н. при 


$j = \ov{1,3}$, где $H_1^j = (H_1^{j1},\dotsc, H_1^{jm})$, 


$j = \ov{1,3}$, 


$H_1^{1j} = \sum \limits_{j=0}^{m_i}\|A_{ij}^1\|


\bigvee \limits_{s=0}^{\displaystyle .}r_{ij}(\cdot,s)$, 


$H_1^{2j} = \sum \limits_{j=0}^{m_i}\|A_{ij}^2\|\,\|H\|_{(\cdot)}


\bigvee \limits_{s=0}^{\displaystyle .}r_{ij}(\cdot,s)$, 


$H_1^{3j} = \sum \limits_{j=0}^{m_i}\|H^{-1}A_{ij}^3\|


\bigvee \limits_{s=0}^{\displaystyle .}r_{ij}(\cdot,s)$, $i = 1,\dotsc, m$, 


$\int \limits_0^{\displaystyle .}(V_2(Hh(0)))(s)dZ(s) \in M_{2p}$ 


для любого $h(0) \in k_{2p}^{n-k}$, то уравнение (1) 


$M_{2p}y$-устойчиво.


\end{cor}





\begin{cor}


Пусть существует $\wh C > 0$ 


такое, что для уравнения (2) имеет место неравенство 


$\int \limits_0^\infty (|H_2^ja^+| + 


\|H_2^jA^+(H_2^j)^\top \|)d\lambda \le \wh C$ п.\,н. при 


$j = \ov{1,3}$, где $H_2^j = (H_2^{j1},\dotsc, H_2^{jm})$, 


$j = \ov{1,3}$, 


$H_2^{1j} = \sum \limits_{j=0}^{m_i}\|\ov A_{ij}^1\|$, 


$H_2^{2j} = \sum \limits_{j=0}^{m_i}\|\ov A_{ij}^2\|\,\|H\|_{(\cdot)}$, 


$H_2^{3j} = \sum \limits_{j=0}^{m_i}\|H^{-1}\ov A_{ij}^3\|$, 


$i = 1,\dotsc, m$, 


$\int \limits_0^{\displaystyle .}(V_2(Hh(0)))(s)dZ(s) \in M_{2p}$ 


для любого $h(0) \in k_{2p}^{n-k}$. Тогда уравнение (2) 


$M_{2p}y$-устойчиво.


\end{cor}





\begin{cor}


Если существует $\wh C > 0$ 


такое, что для уравнения (3) имеет место неравенство 


$\int \limits_0^\infty (|H_3^ja^+| + 


\|H_3^jA^+(H_3^j)^\top \|)d\lambda \le \wh C$ п.\,н. при 


$j = \ov{1,3}$, где $H_3^j = (H_3^{j1},\dotsc, H_3^{jm})$, 


$j = \ov{1,3}$, 


$H_3^{1j} = \|A_i^1\|$, 


$H_3^{2j} = \|A_i^2H\|$, 


$H_3^{3j} = \|H^{-1}A_i^3\|$, 


$i = 1,\dotsc, m$, 


$\int \limits_0^{\displaystyle .}(V_2(Hh(0)))(s)dZ(s) \in M_{2p}$ 


для любого $h(0) \in k_{2p}^{n-k}$, то уравнение (3) 


$M_{2p}y$-устойчиво.


\end{cor}





Отметим, что для получения $M_p^{1/\gamma }y$-устойчивости 


уравнений (1)--(3) можно воспользоваться следствием 1 теоремы 1.





В дальнейшем без ограничения общности предположим, что 


$a^1 = 1$, $A^{11} = 0$, $a^i = 0$ при $i = 2,\dotsc, m$ 


$\lambda^*\times P$ п.\,в. Пусть далее $2p \le q \le \infty$, 


$A^{ij} = 0$\; $\lambda^*\times P$ п.\,в. при $i, j = 1,\dotsc, m$ 


и $i \ne j$, $v(t) = \int \limits_0^t\xi (s)d\lambda (s)$, 


$\Delta v = v(t) - v(s)$, 


$\sup \limits_{t\ge 1}(v(t) - v(t - 1)) < +\infty$,


$\gamma (t) = \exp \{\beta v(t)\}$, $\beta > 0$. 





Для уравнения (1) дополнительно положим


$\|A^l_{1j}\| \le a^l_{1j}$\; $\lambda^*\times P$ п.\,в. и 


$\ov a^l_{1j}=a^l_{1j}\bigvee\limits_{s=0}^{\displaystyle .}


r_{1j}(\cdot,s)\cdot


\xi^{1/q-1} \in L_q^\lambda$ при $l =1, 2$, $j = 0,\dotsc, m_1$, 


$\|A^l_{ij}\|\,|A^{ii}|^{1/2} \le h^l_{ij}$\; $\lambda^*\times P$ п.\,в. и 


$\ov h^l_{ij} = h^l_{ij}\bigvee \limits_{s=0}^{\displaystyle .}r_{ij}(\cdot,s)


\xi^{1/q-1/2} \in L_q^\lambda$ при $l = 1, 2$, 


$i = 2,\dotsc, m$, $j = 0,\dotsc, m_i$, для любого $y \in M_{2p}$ 


имеем $K_1V_3y \in M_{2p}$ и $\|K_1V_3y\|_{M_{2p}} \le \ov c\|y\|_{M_{2p}}$, 


где $\ov c$ --- некоторое положительное число. 


Тогда справедлива





\begin{thm}


Пусть $0 \le l \le m_1$, 


существуют положительные числа $\delta_{1j}$, $j = 0,\dotsc, l$, $\alpha$ 


такие, что $r_{1j}(t,s) = 0$ при $0 \le s \le t - \delta_{1j} < +\infty$, 


$t \in [0, +\infty [$, $j = 0,\dotsc, l$, $\Big\|\sum \limits_{j=0}^l


A_{1j}^1\int \limits_0^{\displaystyle .}d_\nu r_{1j}(\cdot,\nu )


+ \alpha \xi E\Big\| \le g$, 


$\ov g = g\xi^{1/q-1} \in L_q^\lambda $ и выполнено неравенство


\begin{multline*}


\rho \overset{\mathrm{def}}{=}{} (1/\alpha )^{1-1/q}


\bigg[\|\ov g\|_{L_q^\lambda } + 


\sum \limits_{j=0}^l\|\ov a_{1j}^1\|_{L_q^\lambda }


\bigg(\ov \delta_{1j}^{1-1/q}


\sum \limits_{j=0}^{m_1}(\|\ov a_{1j}^1\|_{L_q^\lambda } + 


\ov c\|\ov a_{1j}^2\|_{L_q^\lambda }) + \\


%


+ c_p\ov \delta_{1j}^{1/2-1/q}\sum \limits_{i=2}^m


\sum\limits_{j=0}^{m_i}(\|\ov h_{ij}^1\|_{L_q^\lambda } + 


\ov c\|\ov h_{ij}^2\|_{L_q^\lambda })\bigg) + 


\sum \limits_{j=l+1}^{m_1}\|\ov a_{1j}^1\|_{L_q^\lambda } + 


\ov c\sum \limits_{j=0}^{m_1}\|\ov a_{1j}^2\|_{L_q^\lambda }\bigg] + \\


%


+ c_p(1/2\alpha )^{1/2-1/q}\sum \limits_{i=2}^m


\sum \limits_{j=0}^{m_i}(


\|\ov h_{ij}^1\|_{L_q^\lambda } +


\ov c\|\ov h_{ij}^2\|_{L_q^\lambda }) < 1,


\end{multline*}


где $\ov \delta_{1j} = \sup \limits_{t\ge \delta_{1j}}


\int \limits_{t-\delta_{1j}}^t\xi (s)d\lambda (s)$ при 


$j = 0,\dotsc, l$, и для любого $h(0) \in k_{2p}^{n-k}$ имеем 


\begin{equation}


\int_0^{\displaystyle .}\exp \{-\alpha \Delta v\}(V_2(Hh(0)))(s)dZ(s) 


\in M_{2p}.


\end{equation}


Тогда уравнение $(1)$ $M_{2p}y$-устойчиво. 


Если, кроме того, для некоторого $\beta_0 > 0$ при любых 


$0 \le \beta \le \beta_0$, $z \in M_{2p}$, $h(0) \in k_{2p}^{n-k}$ 


имеем $\gamma K_1V_3(z/\gamma ) \in M_{2p}$, 


\begin{equation}


\int_0^{\displaystyle .}\exp \{-\alpha \Delta v\}\gamma (s)


(V_2(Hh(0)))(s)dZ(s) \in M_{2p}


\end{equation}


и $\|\gamma K_1V_3(z/\gamma )\|_{M_{2p}} \le \ov c_1\|z\|_{M_{2p}}$, 


где $\ov c_1$ --- некоторое положительное число, существуют 


положительные числа 


$\delta_{1j}$, $j = l + 1,\dotsc, m_1$, 


$\delta_{ij}$, $i = 2,\dotsc, m$, $j = 0,\dotsc, m_i$, 


такие, что 


$r_{1j}(t,s) = 0$ при $0 \le s \le t - \delta_{1j} < + \infty$, 


$t \in [0, +\infty [$, $j = l + 1,\dotsc, m_1$, 


$r_{ij}(t,s) = 0$ при $0 \le s \le t - \delta_{ij} < + \infty$, 


$t \in [0, +\infty [$, $i = 2,\dotsc, m$, $j = 0,\dotsc, m_i$, 


то уравнение $(1)$ $M_{2p}^\gamma y$-устойчиво при некотором $\beta > 0$.


\end{thm}





Предположим, что для уравнения (1) выполняются следующие условия: 


$\|A^1_{1j}\| \le a^l_{1j}$, 


$\|A^2_{1j}\|\,\|H\|_{(\cdot)} \le a^l_{1j}$\;


$\lambda^*\times P$ п.\,в. и 


$\ov a^l_{1j} = a^l_{1j}\bigvee \limits_{s=0}^{\displaystyle .}r_{1j}(\cdot,s)


\xi^{1/q-1} \in L_q^\lambda$ при $l =1,2$, $j = 0,\dotsc, m_1$, 


$\|A^1_{ij}\|\,|A^{ii}|^{1/2} \le h^l_{ij}$, 


$\|A^2_{ij}\|\,\|H\|_{(\cdot)}|A^{ii}|^{1/2} \le h^l_{ij}$\;


$\lambda^*\times P$ п.\,в. и 


$\ov h^l_{ij} = h^l_{ij}\bigvee \limits_{s=0}^{\displaystyle .}r_{ij}(\cdot,s)


\xi^{1/q-1/2} \in L_q^\lambda $ при $l = 1, 2$, 


$i = 2,\dotsc, m$, $j = 0,\dotsc, m_i$; для любого $y \in M_{2p}$ 


имеем $\int \limits_0^{\displaystyle .}(\ov V_3y)(s)dZ(s) \in M_{2p}$ и 


$\|\int \limits_0^{\displaystyle .}(\ov V_3y)(s)dZ(s)\|_{M_{2p}}


\le \ov c\|y\|_{M_{2p}}$, 


где $\ov c$ --- некоторое положительное число. 


Тогда справедлива





\begin{thm}


Пусть $0 \le l \le m_1$, 


существуют положительные числа $\delta_{1j}$, $j = 0,\dotsc, l$, $\alpha$ 


такие, что $r_{1j}(t,s) = 0$ при $0 \le s \le t - \delta_{1j} < +\infty$, 


$t \in [0, +\infty [$, $j = 0,\dotsc, l$, $\Big\|\sum \limits_{j=0}^l


A_{1j}^1\int \limits_0^{\displaystyle .}


d_\nu r_{1j}(\cdot,\nu ) + \alpha \xi E\Big\| \le g$, 


$\ov g = g\xi^{1/q-1} \in L_q^\lambda $ и выполнено неравенство


$\rho < 1$, где $\rho $ определен в теореме $2$


и для любого $h(0) \in k_{2p}^{n-k}$ имеем $(6)$.


Тогда уравнение $(1)$ $M_{2p}y$-устойчиво. 


Если, кроме того, для некоторого $\beta_0 > 0$ при любых 


$0 \le \beta \le \beta_0$, $z \in M_{2p}$, $h(0) \in k_{2p}^{n-k}$ 


имеем $(7)$,


$\gamma\int\limits_0^{\displaystyle .}(\ov V_3(z/\gamma))(s)dZ(s)\in M_{2p}$, 


и


$\Big\|\gamma\int\limits_0^{\displaystyle .}


(\ov V_3(z/\gamma ))(s)dZ(s)\Big\|_{M_{2p}} \le 


\ov c_1\|z\|_{M_{2p}}$, 


где $\ov c_1$ --- некоторое положительное число, существуют 


положительные числа 


$\delta_{1j}$, $j = l + 1,\dotsc, m_1$, 


$\delta_{ij}$, $i = 2,\dotsc, m$, $j = 0,\dotsc, m_i$, 


такие, что 


$r_{1j}(t,s) = 0$ при $0 \le s \le t - \delta_{1j} < + \infty$, 


$t \in [0, +\infty [$, $j = l + 1,\dotsc, m_1$, 


$r_{ij}(t,s) = 0$ при $0 \le s \le t - \delta_{ij} < + \infty$, 


$t \in [0, +\infty [$, $i = 2,\dotsc, m$, $j = 0,\dotsc, m_i$, 


то уравнение $(1)$ $M_{2p}^\gamma y$-устойчиво при некотором $\beta > 0$.


\end{thm}





\begin{thebibliography}{9}





\bibitem{1}


Липцер Р.Ш., Ширяев А.Н. {\it Теория мартингалов}. -- М.: 


Наука, 1986. -- 512 с.


\bibitem{2}


Кадиев Р.И., Поносов А.В. {\it Устойчивость линейных стохастических 


функционально-дифференциальных уравнений при постоянно действующих 


возмущениях} // Дифференц. уравнения. --


1992. -- Т.\,28. -- \No\,2. -- С.\,198--207. 


\bibitem{3}


Кадиев Р.И. {\it Существование и единственность решения 


задачи Коши для функционально-дифференциальных уравнений по 


семимартингалу} // Изв. вузов. Математика. --


1995. -- \No\,10. -- С.\,35--40.


\bibitem{4}


Кадиев Р.И. {\it Допустимость пар пространств по части 


переменных для линейных стохастических 


функционально-дифференциальных уравнений} 


// Изв. вузов. Математика. --


1994. -- \No\,4. -- С.\,1--10. \\





\end{thebibliography}





\vskip 2cm





\postup{\it Дагестанский государственный}{$\qquad$\it Поступили}


\postup{\it университет}{\it полный текст $30.06.1994$}


\postup{}{\it краткое сообщение $22.12.1999$}





\label{end}


\end{document}


